\chapter*{ВВЕДЕНИЕ}
\addcontentsline{toc}{chapter}{ВВЕДЕНИЕ}

Растения имеют основополагающее значение для экосистем Земли и жизни человека, играя незаменимую роль в поддержании и улучшении глобального экологического баланса. Методы распознавания образов и обработки изображений используются для автоматизации классификации растений, что позволяет снизить трудоёмкость определения вида и расширить возможности анализа ботанических данных~\cite{ref1101}.

Цель работы~---~разработать метод классификации растений по изображениям листьев с использованием свёрточной нейронной сети.

Для достижения поставленной цели требуется решить следующие задачи:
\begin{itemize}[label=---, itemsep=0pt, parsep=0pt, topsep=0pt]
	\item провести анализ предметной области классификации изображений растений;
	\item формализовать задачу классификации изображений растений;
	\item разработать метод классификации растений по изображениям листьев с использованием свёрточной нейронной сети;
	\item сравнить разработанный метод с существующими методами.
\end{itemize}

\clearpage
